\documentclass[main.tex]{subfiles}
\begin{document}

\chapter{Evaluasi, Refleksi, dan Integrasi Kompetensi}

Bab ini berisi asesmen komprehensif untuk mengukur pencapaian \textbf{CPMK-1} (mengenal algoritma klasik dan modern) dan \textbf{CPMK-2} (memahami konsep dan urgensi di TI).

% ============================================================
% MATERI POKOK
% ============================================================
\input{bab/bab-11/section-11-1}
\input{bab/bab-11/section-11-2}
\input{bab/bab-11/section-11-3}
\input{bab/bab-11/section-11-4}

% ============================================================
% RANGKUMAN PERJALANAN PEMBELAJARAN
% ============================================================

\begin{rangkuman}
Selamat! Anda telah menyelesaikan jalur pembelajaran Kriptografi sesuai silabus OBE.

\textbf{Apa yang Telah Anda Pelajari}:
\begin{itemize}
  \item Konsep, CIA, dan urgensi kriptografi di TI (Sub-CPMK 2.1--2.2)
  \item Matematika pengantar: modulo, prima, invers (2.3)
  \item Cipher klasik dan gagasan frekuensi analisis (1.1)
  \item Stream/block, DES, AES, mode operasi (1.2)
  \item RSA, DH, ElGamal, pengantar ECC (1.3)
  \item Hash, MAC, tanda tangan, PKI/TLS (1.4) dan studi kasus urgensi
\end{itemize}

\textbf{Sumber belajar}: Munir, Stallings, HAC, Schneier, Stinson, Katz, NIST, Boneh (lihat Daftar Pustaka / silabus Bagian~10).
\end{rangkuman}

\begin{transition}
\textbf{Selamat Menyelesaikan Perjalanan Kriptografi!}

Lanjutkan dengan membaca sumber daring silabus (slide Munir ITB, HAC, NIST, Boneh CS255) dan praktik demo ringan. Materi RNG/steganografi tersedia di lampiran opsional.
\end{transition}

\ifSubfilesClassLoaded{
  \renewcommand{\bibname}{Daftar Pustaka}
  \bibliographystyle{plain}
  \bibliography{../references}
}{}

\end{document}
